\begin{lemma}\label{lem:nonic39-positive-witnesses}
Let $K=\mathbb Q(a)$, where
\[
 a^9-5a^8+8a^7-10a^5+4a^4+4a^2-4a+1=0,
\]
and put
\[
 y=\frac{-7a^8+33a^7-46a^6-15a^5+67a^4-7a^3-2a^2-29a+18}{3}.
\]
Define
\[
 r=\frac{a^2y-ay+y-1}{a^2y-a},\qquad
 s=\frac{ay-y+1}{ay},\qquad
 b=rs(r-1),\quad c=s(r-1).
\]
The elliptic curve
\[
 E:\quad Y^2+(1-c)XY-bY=X^3-bX^2
\]
has $E(K)_{\mathrm{tors}}\simeq C_{39}$, generated by $P=(0,0)$.
The diamond translates of the associated closed nonic point on
$X_1(39)$ form four closed nonic divisors, represented by the
physical multipliers $1,2,4,8$.
Their respective multiplier classes modulo sign are
\[
 \{1,16,17\},\quad \{2,5,7\},\quad
 \{4,10,14\},\quad \{8,11,19\}.
\]
At $2$ these four divisors reduce to four distinct noncuspidal closed
points of degree $9$. Their exact Tate-coordinate and Frobenius
bindings are given in \texttt{positive\_divisors39.json}.
\end{lemma}
\begin{proof}
The defining polynomial is irreducible, and its discriminant is
$-3^{14}31^3$. Direct exact arithmetic verifies
$39P=0$, $13P\ne0$ and $3P\ne0$.
At the unramified prime above $29$ defined by $a\equiv10$, one has
$y\equiv20$ and $(b,c)\equiv(17,27)$.
The reduction is
\[
 Y^2+3XY+12Y=X^3+12X^2
\]
over $\mathbb F_{29}$, with discriminant $1$ and exactly $39$ points.
The point count is checked by enumerating all $29^2$ affine pairs.
Good unramified odd-prime reduction injects torsion, so the full
global torsion is $C_{39}$.

For every physical multiplier in
$(\mathbb Z/39\mathbb Z)^*/\{\pm1\}$, the certificate gives the exact
normalized Tate pair for $(E,mP)$. A linear combination of this pair
has an irreducible minimal polynomial of degree $9$, proving that
the marked point has field of definition $K$.
Three explicit automorphisms of $K$ are verified.
There are no others: the automorphism group has order dividing $9$,
and order $9$ would make $K/\mathbb Q$ Galois of odd degree and hence
totally real, contrary to the negative polynomial discriminant.
The action of these three automorphisms on the twelve physical
Tate pairs gives precisely the four triples in the statement.
Since the Tate pair generates $K$, equality of the corresponding
closed moduli points is equivalent to equality under one of these
automorphisms. This proves the global orbit statement.

The defining polynomial modulo $2$ is the irreducible polynomial
$T^9+T^8+1$. All twelve global normalized Tate pairs are integral at
$2$, and their reductions agree exactly with the normalized pairs
of $(\bar E,m\bar P)$ in this fixed $\mathbb F_{512}$.
The three global automorphisms reduce to Frobenius powers $0,3,6$.
Checking all nine Frobenius powers gives exactly the same four
triples, and every pair has joint field of definition
$\mathbb F_{512}$. Thus each global nonic trace reduces to a single
closed nonic point with multiplicity $1$.
The reduced curve is smooth and has $468$ points; its count is also
verified independently by the characteristic-$2$ trace equation.
The normalized Tate pairs are fine level-$39$ structures, so their
Frobenius keys identify the actual modular closed points.
\end{proof}
